% ============================================================
% commun.tex — Définitions partagées entre tous les templates VLM
% (couleurs, commandes, environnements exercice/solution)
% À importer via % ============================================================
% commun.tex — Définitions partagées entre tous les templates VLM
% (couleurs, commandes, environnements exercice/solution)
% À importer via % ============================================================
% commun.tex — Définitions partagées entre tous les templates VLM
% (couleurs, commandes, environnements exercice/solution)
% À importer via % ============================================================
% commun.tex — Définitions partagées entre tous les templates VLM
% (couleurs, commandes, environnements exercice/solution)
% À importer via \input{latex/templates/commun} APRÈS les
% \usepackage nécessaires : xcolor[table], mdframed, verbatim.
% ============================================================

% ── Packages ──
\usepackage[french]{babel}
\usepackage[utf8]{inputenc}
\usepackage[T1]{fontenc}
\usepackage{tgadventor}
\renewcommand*\familydefault{\sfdefault}
\usepackage{amsmath, amssymb, amsthm}
\usepackage[top=2cm, bottom=2cm, left=2cm, right=2cm]{geometry}
\usepackage{enumitem}
\usepackage{multicol}
\usepackage[table]{xcolor}
\usepackage{mdframed}
\usepackage{verbatim}
\usepackage{tabularx}
\usepackage{tkz-euclide}

\definecolor{vlmblue}{RGB}{25, 90, 160}
\definecolor{vlmgray}{RGB}{245, 245, 245}
\definecolor{vlmgreen}{RGB}{0, 130, 80}

\usepackage[thicklines]{cancel}\newcommand\Colxcancel[2][black]{\renewcommand\CancelColor{\color{#1}}\xcancel{#2}}

% ── Commande \sol{question}{réponse} et \colSolution ──
\usepackage{etoolbox}
\providetoggle{question}
\ifdefined\AVECSOLUTIONS
  \togglefalse{question}
\else
  \toggletrue{question}
\fi
\ifdefined\AVECSOLUTIONS
  \colorlet{colSolution}{vlmgreen}
\else
  \colorlet{colSolution}{white}
\fi
% \colSolution utilisable comme commande (\textcolor{\colSolution}{...})
% ou comme couleur directe (\textcolor{colSolution}{...})
\newcommand{\colSolution}{colSolution}
\newcommand{\sol}[2]{\iftoggle{question}{#1}{{\color{colSolution}#2}}}

% ── Réponse mise en couleur ──
\newcommand{\rc}[1]{\textcolor{vlmgreen}{\textbf{#1}}}

% ── Titre de l'exercice dans un encadré ──
\newcommand{\titrexercice}[2]{%
  \ifdefined\AVECSOLUTIONS
    \begin{mdframed}[
      backgroundcolor=vlmgreen!15,
      linecolor=vlmgreen,
      linewidth=1.5pt,
      innerleftmargin=10pt,
      innerrightmargin=10pt,
      innertopmargin=6pt,
      innerbottommargin=6pt,
      skipabove=1em,
      skipbelow=0.5em
    ]
    \textbf{Corrigé #1}%
    \ifx&#2&\else\hfill\textbf{#2}\fi
    \end{mdframed}
  \else
    \begin{mdframed}[
      backgroundcolor=vlmgray,
      linecolor=black,
      linewidth=1.5pt,
      innerleftmargin=10pt,
      innerrightmargin=10pt,
      innertopmargin=6pt,
      innerbottommargin=6pt,
      skipabove=1em,
      skipbelow=0.5em
    ]
    \textbf{Exercice #1}%
    \ifx&#2&\else\hfill\textbf{#2}\fi
    \end{mdframed}
  \fi
}

% ── Corps de l'exercice sans cadre ──
\newenvironment{exercice}[2]{%
  \titrexercice{#1}{#2}
  \begin{list}{}{%
    \setlength{\leftmargin}{0pt}
    \setlength{\rightmargin}{0pt}
    \setlength{\listparindent}{0pt}
    \setlength{\itemindent}{0pt}
    \setlength{\parsep}{0pt}
  }\item[]\noindent\ignorespaces
}{%
  \end{list}
  \vspace{1em}
}

% ── Solution (visible seulement si \AVECSOLUTIONS est défini) ──
\ifdefined\AVECSOLUTIONS
  \newenvironment{solution}[1]{%
    \vspace{0.3em}
    \textbf{\textcolor{vlmgreen}{Solution #1}}
    \vspace{0.3em}\par
  }{%
    \vspace{1em}
  }
\else
  \newenvironment{solution}[1]{\comment}{\endcomment}
\fi
 APRÈS les
% \usepackage nécessaires : xcolor[table], mdframed, verbatim.
% ============================================================

% ── Packages ──
\usepackage[french]{babel}
\usepackage[utf8]{inputenc}
\usepackage[T1]{fontenc}
\usepackage{tgadventor}
\renewcommand*\familydefault{\sfdefault}
\usepackage{amsmath, amssymb, amsthm}
\usepackage[top=2cm, bottom=2cm, left=2cm, right=2cm]{geometry}
\usepackage{enumitem}
\usepackage{multicol}
\usepackage[table]{xcolor}
\usepackage{mdframed}
\usepackage{verbatim}
\usepackage{tabularx}
\usepackage{tkz-euclide}

\definecolor{vlmblue}{RGB}{25, 90, 160}
\definecolor{vlmgray}{RGB}{245, 245, 245}
\definecolor{vlmgreen}{RGB}{0, 130, 80}

\usepackage[thicklines]{cancel}\newcommand\Colxcancel[2][black]{\renewcommand\CancelColor{\color{#1}}\xcancel{#2}}

% ── Commande \sol{question}{réponse} et \colSolution ──
\usepackage{etoolbox}
\providetoggle{question}
\ifdefined\AVECSOLUTIONS
  \togglefalse{question}
\else
  \toggletrue{question}
\fi
\ifdefined\AVECSOLUTIONS
  \colorlet{colSolution}{vlmgreen}
\else
  \colorlet{colSolution}{white}
\fi
% \colSolution utilisable comme commande (\textcolor{\colSolution}{...})
% ou comme couleur directe (\textcolor{colSolution}{...})
\newcommand{\colSolution}{colSolution}
\newcommand{\sol}[2]{\iftoggle{question}{#1}{{\color{colSolution}#2}}}

% ── Réponse mise en couleur ──
\newcommand{\rc}[1]{\textcolor{vlmgreen}{\textbf{#1}}}

% ── Titre de l'exercice dans un encadré ──
\newcommand{\titrexercice}[2]{%
  \ifdefined\AVECSOLUTIONS
    \begin{mdframed}[
      backgroundcolor=vlmgreen!15,
      linecolor=vlmgreen,
      linewidth=1.5pt,
      innerleftmargin=10pt,
      innerrightmargin=10pt,
      innertopmargin=6pt,
      innerbottommargin=6pt,
      skipabove=1em,
      skipbelow=0.5em
    ]
    \textbf{Corrigé #1}%
    \ifx&#2&\else\hfill\textbf{#2}\fi
    \end{mdframed}
  \else
    \begin{mdframed}[
      backgroundcolor=vlmgray,
      linecolor=black,
      linewidth=1.5pt,
      innerleftmargin=10pt,
      innerrightmargin=10pt,
      innertopmargin=6pt,
      innerbottommargin=6pt,
      skipabove=1em,
      skipbelow=0.5em
    ]
    \textbf{Exercice #1}%
    \ifx&#2&\else\hfill\textbf{#2}\fi
    \end{mdframed}
  \fi
}

% ── Corps de l'exercice sans cadre ──
\newenvironment{exercice}[2]{%
  \titrexercice{#1}{#2}
  \begin{list}{}{%
    \setlength{\leftmargin}{0pt}
    \setlength{\rightmargin}{0pt}
    \setlength{\listparindent}{0pt}
    \setlength{\itemindent}{0pt}
    \setlength{\parsep}{0pt}
  }\item[]\noindent\ignorespaces
}{%
  \end{list}
  \vspace{1em}
}

% ── Solution (visible seulement si \AVECSOLUTIONS est défini) ──
\ifdefined\AVECSOLUTIONS
  \newenvironment{solution}[1]{%
    \vspace{0.3em}
    \textbf{\textcolor{vlmgreen}{Solution #1}}
    \vspace{0.3em}\par
  }{%
    \vspace{1em}
  }
\else
  \newenvironment{solution}[1]{\comment}{\endcomment}
\fi
 APRÈS les
% \usepackage nécessaires : xcolor[table], mdframed, verbatim.
% ============================================================

% ── Packages ──
\usepackage[french]{babel}
\usepackage[utf8]{inputenc}
\usepackage[T1]{fontenc}
\usepackage{tgadventor}
\renewcommand*\familydefault{\sfdefault}
\usepackage{amsmath, amssymb, amsthm}
\usepackage[top=2cm, bottom=2cm, left=2cm, right=2cm]{geometry}
\usepackage{enumitem}
\usepackage{multicol}
\usepackage[table]{xcolor}
\usepackage{mdframed}
\usepackage{verbatim}
\usepackage{tabularx}
\usepackage{tkz-euclide}

\definecolor{vlmblue}{RGB}{25, 90, 160}
\definecolor{vlmgray}{RGB}{245, 245, 245}
\definecolor{vlmgreen}{RGB}{0, 130, 80}

\usepackage[thicklines]{cancel}\newcommand\Colxcancel[2][black]{\renewcommand\CancelColor{\color{#1}}\xcancel{#2}}

% ── Commande \sol{question}{réponse} et \colSolution ──
\usepackage{etoolbox}
\providetoggle{question}
\ifdefined\AVECSOLUTIONS
  \togglefalse{question}
\else
  \toggletrue{question}
\fi
\ifdefined\AVECSOLUTIONS
  \colorlet{colSolution}{vlmgreen}
\else
  \colorlet{colSolution}{white}
\fi
% \colSolution utilisable comme commande (\textcolor{\colSolution}{...})
% ou comme couleur directe (\textcolor{colSolution}{...})
\newcommand{\colSolution}{colSolution}
\newcommand{\sol}[2]{\iftoggle{question}{#1}{{\color{colSolution}#2}}}

% ── Réponse mise en couleur ──
\newcommand{\rc}[1]{\textcolor{vlmgreen}{\textbf{#1}}}

% ── Titre de l'exercice dans un encadré ──
\newcommand{\titrexercice}[2]{%
  \ifdefined\AVECSOLUTIONS
    \begin{mdframed}[
      backgroundcolor=vlmgreen!15,
      linecolor=vlmgreen,
      linewidth=1.5pt,
      innerleftmargin=10pt,
      innerrightmargin=10pt,
      innertopmargin=6pt,
      innerbottommargin=6pt,
      skipabove=1em,
      skipbelow=0.5em
    ]
    \textbf{Corrigé #1}%
    \ifx&#2&\else\hfill\textbf{#2}\fi
    \end{mdframed}
  \else
    \begin{mdframed}[
      backgroundcolor=vlmgray,
      linecolor=black,
      linewidth=1.5pt,
      innerleftmargin=10pt,
      innerrightmargin=10pt,
      innertopmargin=6pt,
      innerbottommargin=6pt,
      skipabove=1em,
      skipbelow=0.5em
    ]
    \textbf{Exercice #1}%
    \ifx&#2&\else\hfill\textbf{#2}\fi
    \end{mdframed}
  \fi
}

% ── Corps de l'exercice sans cadre ──
\newenvironment{exercice}[2]{%
  \titrexercice{#1}{#2}
  \begin{list}{}{%
    \setlength{\leftmargin}{0pt}
    \setlength{\rightmargin}{0pt}
    \setlength{\listparindent}{0pt}
    \setlength{\itemindent}{0pt}
    \setlength{\parsep}{0pt}
  }\item[]\noindent\ignorespaces
}{%
  \end{list}
  \vspace{1em}
}

% ── Solution (visible seulement si \AVECSOLUTIONS est défini) ──
\ifdefined\AVECSOLUTIONS
  \newenvironment{solution}[1]{%
    \vspace{0.3em}
    \textbf{\textcolor{vlmgreen}{Solution #1}}
    \vspace{0.3em}\par
  }{%
    \vspace{1em}
  }
\else
  \newenvironment{solution}[1]{\comment}{\endcomment}
\fi
 APRÈS les
% \usepackage nécessaires : xcolor[table], mdframed, verbatim.
% ============================================================

% ── Packages ──
\usepackage[french]{babel}
\usepackage[utf8]{inputenc}
\usepackage[T1]{fontenc}
\usepackage{tgadventor}
\renewcommand*\familydefault{\sfdefault}
\usepackage{amsmath, amssymb, amsthm}
\usepackage[top=2cm, bottom=2cm, left=2cm, right=2cm]{geometry}
\usepackage{enumitem}
\usepackage{multicol}
\usepackage[table]{xcolor}
\usepackage{mdframed}
\usepackage{verbatim}
\usepackage{tabularx}
\usepackage{tkz-euclide}

\definecolor{vlmblue}{RGB}{25, 90, 160}
\definecolor{vlmgray}{RGB}{245, 245, 245}
\definecolor{vlmgreen}{RGB}{0, 130, 80}

\usepackage[thicklines]{cancel}\newcommand\Colxcancel[2][black]{\renewcommand\CancelColor{\color{#1}}\xcancel{#2}}

% ── Commande \sol{question}{réponse} et \colSolution ──
\usepackage{etoolbox}
\providetoggle{question}
\ifdefined\AVECSOLUTIONS
  \togglefalse{question}
\else
  \toggletrue{question}
\fi
\ifdefined\AVECSOLUTIONS
  \colorlet{colSolution}{vlmgreen}
\else
  \colorlet{colSolution}{white}
\fi
% \colSolution utilisable comme commande (\textcolor{\colSolution}{...})
% ou comme couleur directe (\textcolor{colSolution}{...})
\newcommand{\colSolution}{colSolution}
\newcommand{\sol}[2]{\iftoggle{question}{#1}{{\color{colSolution}#2}}}

% ── Réponse mise en couleur ──
\newcommand{\rc}[1]{\textcolor{vlmgreen}{\textbf{#1}}}

% ── Titre de l'exercice dans un encadré ──
\newcommand{\titrexercice}[2]{%
  \ifdefined\AVECSOLUTIONS
    \begin{mdframed}[
      backgroundcolor=vlmgreen!15,
      linecolor=vlmgreen,
      linewidth=1.5pt,
      innerleftmargin=10pt,
      innerrightmargin=10pt,
      innertopmargin=6pt,
      innerbottommargin=6pt,
      skipabove=1em,
      skipbelow=0.5em
    ]
    \textbf{Corrigé #1}%
    \ifx&#2&\else\hfill\textbf{#2}\fi
    \end{mdframed}
  \else
    \begin{mdframed}[
      backgroundcolor=vlmgray,
      linecolor=black,
      linewidth=1.5pt,
      innerleftmargin=10pt,
      innerrightmargin=10pt,
      innertopmargin=6pt,
      innerbottommargin=6pt,
      skipabove=1em,
      skipbelow=0.5em
    ]
    \textbf{Exercice #1}%
    \ifx&#2&\else\hfill\textbf{#2}\fi
    \end{mdframed}
  \fi
}

% ── Corps de l'exercice sans cadre ──
\newenvironment{exercice}[2]{%
  \titrexercice{#1}{#2}
  \begin{list}{}{%
    \setlength{\leftmargin}{0pt}
    \setlength{\rightmargin}{0pt}
    \setlength{\listparindent}{0pt}
    \setlength{\itemindent}{0pt}
    \setlength{\parsep}{0pt}
  }\item[]\noindent\ignorespaces
}{%
  \end{list}
  \vspace{1em}
}

% ── Solution (visible seulement si \AVECSOLUTIONS est défini) ──
\ifdefined\AVECSOLUTIONS
  \newenvironment{solution}[1]{%
    \vspace{0.3em}
    \textbf{\textcolor{vlmgreen}{Solution #1}}
    \vspace{0.3em}\par
  }{%
    \vspace{1em}
  }
\else
  \newenvironment{solution}[1]{\comment}{\endcomment}
\fi
